\subsubsection*{2.2.6.}
我们将给出 D 和 E 中的圈 $\delta$ 与 $\delta'$ 在平展拓扑下的局部描述.

$n$ 为偶数时很容易. 若 $2\nmid k$ ($\Lambda=\mathbb Z/k$), 则 $\pm\delta$ 由 (2.2.5.3) 和 (2.2.5.4) 刻画, 即由
\begin{equation}\tag{2.2.6.1}
(\delta,\delta)=(-1)^m\,2.
\end{equation}
刻画. 一般地, $\pm\delta$ 是任意一个以 $\mathbb Z/2^ak$ 为系数的上同调中满足 (2.2.6.1) 的圈 $\widetilde\delta$ 的像, 其中 $a$ 充分大.

假设 $n$ 为奇数. 则 $x_0$ 为有理点(2.2.2). 设 $X_{s(0)}$ 为 $X_s$ 在 $x_0$ 处的 Hensel 化, $\widetilde X_{s(0)}$ 为其在 $x_0$ 处的爆破, $Y_0$ 为 $x_0$ 在 $\widetilde X_{s(0)}$ 中的逆像(一个二次超曲面), 并置
\[
X^*_{s(0)}=X_{s(0)}-\{x_0\}
=\widetilde X_{s(0)}-Y_0.
\]
考虑复合映射(参见 XIII, 2.4.6.1)
\begin{equation}\tag{2.2.6.2}
\begin{tikzcd}[column sep=huge,row sep=huge]
H^{n-1}(Y_0,\Lambda(m))
  & H^{n-1}(\widetilde X_{s(0)},\Lambda(m))
      \arrow[l,"\sim"'] \arrow[r]
  & H^{n-1}(X^*_{s(0)},\Lambda(m)) \arrow[d] \\
{}
  & H^n_{\{x_0\}}\bigl(X_s,R\Psi_{\bar\eta}(\Lambda(m))\bigr)
  & H^n_{\{x_0\}}(X_s,\Lambda(m)) \arrow[l]
\end{tikzcd}
\end{equation}

\subsubsection*{引理 2.2.7.}
\underline{映射 (2.2.6.2) 将 $H^n_{\{x_0\}}(X_s,R\Psi_{\bar\eta}(\Lambda(m)))$ 等同于 $2m$ 维二次超曲面 $Y_0$ 的中间维数上同调的本原商; 两个类 $\pm\delta$ 是此本原商的自然生成元的像.}

圈 $\pm\delta$ 由其在 $H_c^n(X_{\bar\eta},\Lambda(m))$ 中的像刻画. 因此利用 (2.1.7.1), 只需证明复合映射
\begin{center}
\begin{tikzcd}[column sep=huge,row sep=huge]
H^{n-1}(Y_0,\Lambda(m))
  & H^{n-1}(\widetilde X_s,\Lambda(m))
      \arrow[l,"\sim"'] \arrow[r]
  & H^{n-1}(X_s^*,\Lambda(m)) \arrow[d] \\
H_c^n(X_{\bar\eta},\Lambda(m))
  & H_c^n(X_s,\Lambda(m)) \arrow[l,"\operatorname{sp}"']
  & H^n_{\{x_0\}}(X_s,\Lambda(m)) \arrow[l]
\end{tikzcd}
\end{center}
将第一个群的本原商的特定生成元映到最后一个群的特定生成元.

应用 2.1.8. 在相差一个符号的意义下, 此映射可改写为
\begin{center}
\begin{tikzcd}[column sep=huge,row sep=huge]
H^{n-1}(Y_s,\Lambda(m)) \arrow[r,"\partial"] \arrow[d,equal]
  & H_c^n(X_s,\Lambda(m)) \arrow[d] \\
H^{n-1}(Y_{\bar\eta},\Lambda(m)) \arrow[r]
  & H_c^n(X_{\bar\eta},\Lambda(m)),
\end{tikzcd}
\end{center}
从而得到 2.2.7.
\footnote{在此证明中, 印本将刻画次数为 $n$ 的圈 $\delta$ 的群写为 $H_c^i$; 所列映射及全部上下文均要求 $H_c^n$. 印本还将 (2.2.6.2) 中的支集点简写为 $x$; 英译使用固定的奇点 $x_0$.}